% ============================================================
%  GUÍA DE ESTUDIO — HISTORIA CONTEMPORÁNEA (4.º AÑO)
%  Guía 2: Gobierno de Juan Vicente Gómez (1908–1935)
%  Reedición de estudio de la guía del Prof. José Rambert
%  (reestructurada con su permiso, en sesiones Pomodoro)
%  Centro Educativo Las Américas — Maturín, Edo. Monagas
%  Compilar:  latexmk -xelatex guia_1_2_gomez.tex
% ============================================================

\documentclass[12pt, letterpaper]{article}

\usepackage{mathwizards-guia}
\usepackage{tabularx}
\usepackage{longtable}
\usepackage{needspace}

% Evitar títulos de sesión huérfanos al final de la página:
% si queda menos de 7 cm, la sesión empieza en página nueva.
\makeatletter
\let\mw@oldsection\section
\newcommand{\mw@prepage}{%
  \par
  \ifdim\pagegoal=\maxdimen\else
    \ifdim\dimexpr\pagegoal-\pagetotal\relax<7cm\relax
      \newpage
    \fi
  \fi}
\renewcommand{\section}{\mw@prepage\mw@oldsection}
\makeatother

% ── Comandos propios de esta guía ────────────────────────────

% Etiquetas de nivel de dificultad
\newcommand{\nUno}{\textcolor{mwnaranja}{\textbf{N1}}}
\newcommand{\nDos}{\textcolor{mwrojonaranja}{\textbf{N2}}}
\newcommand{\nTres}{\textcolor{mwmagenta}{\textbf{N3}}}
\newcommand{\nCuatro}{\textcolor{mwvioleta}{\textbf{N4}}}

% Cabecera de sesión (meta + tiempo)
\newmdenv[
  backgroundcolor = mwgrisclaro,
  linecolor       = mwprimario,
  linewidth       = 2pt,
  topline         = false,
  bottomline      = false,
  rightline       = false,
  leftline        = true,
  innertopmargin  = 6pt,
  innerbottommargin = 6pt,
  innerleftmargin = 10pt,
  innerrightmargin = 8pt,
  skipabove       = 6pt,
  skipbelow       = 8pt
]{cabecerasesion}

% Caja de descanso
\newmdenv[
  backgroundcolor = mwnaranja!10,
  linecolor       = mwnaranja,
  linewidth       = 1pt,
  roundcorner     = 4pt,
  innertopmargin  = 5pt,
  innerbottommargin = 5pt,
  innerleftmargin = 10pt,
  innerrightmargin = 10pt,
  skipabove       = 8pt,
  skipbelow       = 6pt
]{descanso}

% Preguntas de análisis
\newlist{preguntas}{enumerate}{1}
\setlist[preguntas]{
  label     = \textbf{\arabic*.},
  leftmargin= *,
  itemsep   = 4pt,
  topsep    = 4pt
}

% Ficha de sesión
\newcommand{\metasesion}[2]{%
  \begin{cabecerasesion}
    \small
    \textbf{\color{mwprimario}Tiempo:} #1 \quad$\cdot$\quad
    \textbf{\color{mwprimario}Meta:} #2
  \end{cabecerasesion}%
}

\newcommand{\pausa}[1]{%
  \begin{descanso}
    \small\textbf{\color{mwnaranja}Descanso (5 min).} #1
  \end{descanso}%
}

\mwguia{Guía 2 — Juan Vicente Gómez (1908--1935)}%
       {Prof. José Rambert $\cdot$ 4.º Año}

\begin{document}
% ============================================================

% ── Portada ─────────────────────────────────────────────────
\mwportada{Guía 2 — Juan Vicente Gómez}
          {Historia Contemporánea de Venezuela (1908--1935)}
          {10 sesiones Pomodoro para estudiar con método}

\vspace{2pt}
\begin{cuadroinfo}
  \small
  \textbf{Institución:} Centro Educativo Las Américas --- Maturín, Edo. Monagas
  \quad$\cdot$\quad \textbf{Secciones:} B, C, D \quad$\cdot$\quad
  \textbf{Fecha:} 16/10/26 \\[3pt]
  \textbf{Docente:} Prof. José Rambert \\[3pt]
  \textbf{Objetivo 2.1:} Analizar el proceso político vivido en Venezuela con la
  llegada de los Andinos al Poder entre 1908 y 1935 con el gobierno de
  Juan Vicente Gómez.
\end{cuadroinfo}

\vspace{6pt}

\begin{cuadroadvertencia}
  \small
  \textbf{Créditos y honestidad intelectual.} Esta guía es una
  \textbf{reedición de estudio} de la guía original del
  \textbf{Prof. José Rambert}. No se altera el
  contenido histórico: se \textbf{reorganiza} la información en sesiones de
  25/5 minutos (Pomodoro) y se añaden \textbf{preguntas de análisis} al final de
  cada sesión. Si algún dato de esta versión no coincide con lo explicado en
  clase, \textbf{prevalece la versión del profesor}.
\end{cuadroadvertencia}

\begin{cuadroteoria}
  \small
  \textbf{Cómo usar esta guía.}
  \begin{itemize}[leftmargin=*, itemsep=2pt, topsep=2pt]
    \item Cada sesión dura \textbf{25 min de trabajo + 5 min de descanso}.
    \item Dentro de los 25 min: $\approx$ 16--18 min de lectura y
          $\approx$ 7--9 min para responder.
    \item Cada sesión deja un \textbf{producto} (tabla, línea de tiempo,
          esquema o párrafo). No avances sin escribirlo.
    \item Después de \textbf{4 sesiones}, toma un descanso largo de 15--20 min.
    \item Materiales: cuaderno, lápiz, cronómetro y esta guía.
  \end{itemize}
  \vspace{2pt}
  \textbf{Niveles de las preguntas:}\quad
  \nUno{} recordar / localizar \quad
  \nDos{} comprender / explicar \quad
  \nTres{} analizar / comparar \quad
  \nCuatro{} evaluar / argumentar (reto)
  \\[3pt]
  \textit{Regla de oro:} si no puedes responder una pregunta \nUno{}, vuelve al
  texto antes de intentar la \nCuatro{}. El análisis sin datos es opinión.
\end{cuadroteoria}

\vspace{4pt}

\begin{center}
\small
\begin{tabular}{@{}c p{12.4cm} r@{}}
  \toprule
  \textbf{\#} & \textbf{Sesión} & \textbf{Tiempo} \\
  \midrule
  1  & El escenario: de Cipriano Castro a Gómez (1899--1908) & 30 min \\
  2  & El hombre: biografía de Juan Vicente Gómez            & 30 min \\
  3  & El sistema: cómo se sostenía el poder                 & 30 min \\
  4  & La ley como herramienta: períodos y reelección        & 30 min \\
  5  & La fuerza: ejército, represión y nepotismo            & 30 min \\
  6  & La economía I: la Venezuela agrícola (1908--1921)     & 30 min \\
  7  & La economía II: petróleo, crisis de 1929 y deuda      & 30 min \\
  8  & La sociedad: población, éxodo y estratos              & 30 min \\
  9  & La cultura: educación, positivismo y prensa           & 30 min \\
  10 & La oposición y la Generación del 28                   & 30 min \\
  11 & Cierre: síntesis, cronología y simulacro              & 30 min \\
  \bottomrule
\end{tabular}
\end{center}

\noindent\small\textbf{Tiempo total:} $\approx$ 5 h 30 min, incluidos los
descansos.

\newpage

% ============================================================
\section{El escenario: de Castro a Gómez (1899--1908)}
% ============================================================

\metasesion{25 + 5 min}{explicar cómo la carrera militar de Gómez y la ausencia
de Castro crearon la oportunidad política de 1908.}

Antes de estudiar a Gómez hay que entender el terreno:

\begin{itemize}[leftmargin=*, itemsep=3pt]
  \item \textbf{1899:} la \textbf{Revolución Liberal Restauradora} lleva a
        \textbf{Cipriano Castro} al poder. Gómez participa y queda dentro del
        nuevo grupo gobernante.
  \item \textbf{1902:} frente a la \textbf{Revolución Libertadora}, Gómez
        obtiene la victoria y gana \textbf{prestigio militar}. Ese prestigio es
        su capital político.
  \item \textbf{1908:} Castro viaja \textbf{enfermo a Europa} y le deja el
        poder. Desde el \textbf{19 de diciembre de 1908} y hasta el día de su
        muerte, Gómez gobierna de forma dictatorial.
\end{itemize}

\begin{cuadropasos}
  \small
  \textbf{Producto de la sesión.} Línea de tiempo con cuatro hitos
  (1899 -- 1902 -- 1908 -- 1935) y una frase por hito.
\end{cuadropasos}

\subsection*{Preguntas de análisis}
\begin{preguntas}
  \item \nUno{} ¿Qué dos hechos militares consolidaron la carrera de Gómez antes
        de 1908?
  \item \nUno{} ¿Por qué Gómez queda al mando de la presidencia en diciembre de
        1908?
  \item \nDos{} Explica con tus palabras la diferencia entre ser
        \textbf{presidente provisional} y ser \textbf{presidente
        constitucional} en el caso de Gómez.
  \item \nTres{} ¿La enfermedad de Castro fue solo una casualidad? Analiza qué
        condiciones tenía ya Gómez preparadas antes de 1908.
  \item \nCuatro{} ¿Se puede afirmar que sin Castro no habría existido Gómez
        como dictador? Responde usando al menos \textbf{dos datos} de esta
        sesión.
\end{preguntas}

\pausa{Antes de continuar, di en voz alta la respuesta a la pregunta 1.}

% ============================================================
\section{El hombre: biografía de Juan Vicente Gómez}
% ============================================================

\metasesion{25 + 5 min}{reconstruir la trayectoria de Gómez y reconocer cómo su
biografía anticipa su forma de gobernar.}

{\small
\begin{longtable}{@{}>{\raggedright\arraybackslash}p{3.3cm} >{\raggedright\arraybackslash}p{12.3cm}@{}}
  \toprule
  \textbf{Etapa} & \textbf{Datos clave} \\
  \midrule
  \endfirsthead
  \toprule
  \textbf{Etapa} & \textbf{Datos clave} \\
  \midrule
  \endhead

  \textbf{Origen} &
  Nace el \textbf{24 de julio de 1857} en la hacienda \textbf{La Mulera}, cerca
  de San Antonio del Táchira. Hijo de \textbf{Pedro Cornelio Gómez} y
  \textbf{Hermenegilda Chacón Alarcón}; tuvo \textbf{12 hermanos} más. \\

  \textbf{Familia y tierra} &
  Al morir su padre (\textbf{14 de enero de 1883}) se convierte en
  \guillemotleft cabeza de familia\guillemotright{} y se dedica al campo y a la
  cría de ganado. \\

  \textbf{Política regional} &
  En \textbf{1888} se hace amigo de \textbf{Cipriano Castro}, gobernador del
  Táchira. \\

  \textbf{Guerra} &
  En \textbf{1892}, como comisario de Guerra y con rango de coronel, defiende al
  gobierno de \textbf{Raimundo Andueza Palacio}, derrocado por la
  \guillemotleft Revolución Legalista\guillemotright. Pasa \textbf{siete años
  exiliado} cerca de Cúcuta. \\

  \textbf{Poder} &
  Apoya la Revolución Liberal Restauradora (\textbf{1899}). Vence en la
  Revolución Libertadora (\textbf{1902}). \\

  \textbf{Dictadura} &
  Gobierna del \textbf{19/12/1908} hasta su muerte, durante
  \textbf{tres mandatos presidenciales}, con varias reformas constitucionales
  para alargar el período y permitir la reelección. \\

  \textbf{Presidencias} &
  Presidente Provisional (1908); luego constitucional en
  \textbf{1910--1915}, \textbf{1915--1922}, \textbf{1922--1929} y
  \textbf{1929--1936}. Para el último período fingió no aceptar y recomendó a
  \textbf{Juan Bautista Pérez}; en \textbf{1931} lo hizo renunciar y asumió
  hasta su muerte. \\

  \textbf{Ausencias del cargo} &
  Se separó varias veces de la Presidencia, pero conservó la
  \textbf{Comandancia en Jefe del Ejército}: \textbf{agosto 1913--enero 1914}
  (Dr. José Gil Fortoul) y \textbf{1915--1922} (Dr. Victoriano Márquez
  Bustillos). \\

  \textbf{Riqueza} &
  Mayor propietario territorial: hatos de \textbf{centenares de miles de
  hectáreas}, principal proveedor de carne del país, \textbf{más del 60\,\%} del
  ganado nacional. Monopolizó el \textbf{juego} (legalizado), la
  \textbf{navegación fluvial y costanera}, la \textbf{electricidad} en varias
  ciudades, \textbf{ingenios de azúcar}, haciendas de \textbf{café y cacao} y
  pequeñas industrias (telas, jabones, cigarrillos, velas, vidrio, aceites,
  hoteles). \\

  \textbf{Otros hechos} &
  En \textbf{1930} cancela la deuda externa como homenaje a Bolívar en el
  centenario de su muerte. Inhabilitó los partidos de oposición y castigó
  duramente la delincuencia. Hermanos, hijos y familiares ocuparon altos
  cargos. \\

  \textbf{Muerte} &
  \textbf{18 de diciembre de 1935}, en \textbf{Maracay}. \\
  \bottomrule
\end{longtable}
}

\begin{cuadropasos}
  \small
  \textbf{Producto de la sesión.} Ficha biográfica de 8 líneas que contenga las
  palabras: \textit{Táchira -- Castro -- 1902 -- 1908 -- Ejército -- ganado --
  Maracay}.
\end{cuadropasos}

\subsection*{Preguntas de análisis}
\begin{preguntas}
  \item \nUno{} Localiza fecha y lugar de nacimiento, y fecha y lugar de muerte.
  \item \nUno{} Enumera \textbf{tres} negocios que Gómez monopolizó.
  \item \nDos{} ¿Cómo pasa de \guillemotleft cabeza de familia\guillemotright{} a
        caudillo nacional? Ordena: 1883, 1888, 1892, 1899, 1902.
  \item \nTres{} ¿Por qué controlar el ganado y la venta de carne era también
        una forma de poder \textbf{político} y no solo económico?
  \item \nCuatro{} El texto afirma que Gómez
        \guillemotleft superó a los caudillos que le precedieron\guillemotright.
        ¿En qué aspectos la superación es real y en cuáles es solo una dictadura
        más eficiente? Argumenta.
\end{preguntas}

\pausa{Bebe agua y repite mentalmente las cuatro fechas de la pregunta 3.}

% ============================================================
\section{El sistema: cómo se sostenía el poder}
% ============================================================

\metasesion{25 + 5 min}{identificar los mecanismos políticos que hicieron
posible 27 años de gobierno.}

\begin{description}[leftmargin=1.2em, itemsep=5pt, style=nextline]
  \item[\textbf{a) La duración.}]
  Gobernó \textbf{27 años (1908--1935)}, unas veces como presidente directo y
  otras colocando en el poder, de forma nominal, a un amigo suyo.

  \item[\textbf{b) El Congreso.}]
  Estuvo \textbf{siempre a su servicio} y reformó la Constitución en varias
  ocasiones con \textbf{dos objetivos básicos}:
  \begin{enumerate}[leftmargin=*, itemsep=1pt, topsep=2pt]
    \item \textbf{variar el período presidencial} y \textbf{permitir la
          reelección};
    \item dar a la \textbf{Comandancia de las Fuerzas Armadas} el rango de
          \textbf{co-presidencia} / jefatura separada. Así, en algunas
          ocasiones Gómez era el Comandante de las Fuerzas Armadas y en otras
          tenía ese cargo junto con el de presidente.
  \end{enumerate}

  \item[\textbf{c) El Consejo de Gobierno.}]
  Lo mantuvo y después lo eliminó. En él estuvieron inicialmente
  \textbf{generales amarillos y nacionalistas}; luego desistió de ellos y anuló
  su acción, igual que a los \textbf{conservadores azules}.
  \textbf{No hubo partidos políticos} durante el gomecismo.

  \item[\textbf{d) La infraestructura del control.}]
  Organizó la defensa con un \textbf{ejército modernizado y bien equipado}
  (reforma iniciada por Castro). Desarrolló una \textbf{red de carreteras de
  macadán} ---algunas de asfalto--- que unía las ciudades con las
  \textbf{plazas militares}.

  \item[\textbf{e) Represión legal e ilegal.}]
  La dictadura se manifestó en el \textbf{manejo del Congreso} y en la
  \textbf{mano dura}: \textbf{cadenas de exilio}, cárceles
  (\textbf{La Rotunda} en Caracas, \textbf{Las Tres Torres} en Barquisimeto,
  \textbf{Castillo de Puerto Cabello}), \textbf{trabajos forzados} en la
  construcción de carreteras y \textbf{torturas} frecuentes.

  \item[\textbf{f) Las alianzas.}]
  Se puso del lado de las llamadas \textbf{«fuerzas vivas»} (la burguesía
  capitalista), tuvo el \textbf{apoyo constante del gobierno norteamericano} y
  de otros países, y aplicó una política de \textbf{grandes concesiones}, sobre
  todo petroleras, para obtener respaldo extranjero.
\end{description}

\begin{cuadropasos}
  \small
  \textbf{Producto de la sesión.} Esquema de cinco flechas:
  \textit{Congreso $\rightarrow$ Constitución $\rightarrow$ reelección
  $\rightarrow$ Ejército $\rightarrow$ concesiones}.
\end{cuadropasos}

\subsection*{Preguntas de análisis}
\begin{preguntas}
  \item \nUno{} ¿Cuántos años gobernó y cuántos partidos políticos existieron en
        ese período?
  \item \nUno{} ¿Qué era el Consejo de Gobierno y qué ocurrió con él?
  \item \nDos{} Explica los \textbf{dos objetivos} de las reformas
        constitucionales.
  \item \nTres{} ¿Cómo un Congreso \guillemotleft a su servicio\guillemotright{}
        le da \textbf{apariencia de legalidad} a la dictadura?
  \item \nCuatro{} ¿Qué diferencia hay entre una dictadura que
        \textbf{elimina} las instituciones y otra que las \textbf{vacía por
        dentro}? Aplica el concepto al Congreso y al Consejo de Gobierno
        gomecistas.
\end{preguntas}

\pausa{Compara tu esquema con el de un compañero o repítelo en voz alta.}

% ============================================================
\section{La ley como herramienta: períodos y reelección}
% ============================================================

\metasesion{25 + 5 min}{leer la cronología como un mecanismo: cada reforma
constitucional resuelve un problema de poder.}

{\small
\begin{longtable}{@{}>{\raggedright\arraybackslash}p{2.6cm}
                       >{\raggedright\arraybackslash}p{5.2cm}
                       >{\raggedright\arraybackslash}p{7.8cm}@{}}
  \toprule
  \textbf{Período} & \textbf{Gobernante} & \textbf{Reforma y qué cambió} \\
  \midrule
  \endfirsthead
  \toprule
  \textbf{Período} & \textbf{Gobernante} & \textbf{Reforma y qué cambió} \\
  \midrule
  \endhead

  \textbf{1908--1913} & Gral. Juan Vicente Gómez (1.ª vez) &
  \textbf{1909:} reduce el período presidencial a \textbf{4 años} (también el de
  senadores y diputados). Crea el \textbf{Consejo de Gobierno}, que elige a los
  miembros del Congreso. \\

  \textbf{1913} & Dr. José Gil Fortoul (encargado por Gómez) & --- \\

  \textbf{1914} & Dr. Victoriano Márquez Bustillo (encargado) &
  Aumenta el período presidencial a \textbf{7 años} y \textbf{permite la
  reelección}. Las \textbf{faltas temporales} las suple el ministro que el
  presidente elija; las \textbf{faltas absolutas}, el presidente de la
  \textbf{Alta Corte Federal y de Casación}, que convoca al Congreso a nueva
  elección. \\

  \textbf{1914--1922} & Gral. Gómez (2.ª vez) &
  Crea los cargos de \textbf{Primer y Segundo Vicepresidente}.
  \textbf{Elimina} el cargo de Comandante General del Ejército. \\

  \textbf{1922--1929} & Gral. Gómez (3.ª vez) &
  \textbf{1925:} el Poder Ejecutivo se sitúa \textbf{en cualquier estado donde
  esté el presidente}. \textbf{Despoja a los estados} del derecho a elegir a sus
  gobernadores. \\

  \textbf{1929--1931} & Dr. Juan Bautista Pérez (puesto por Gómez; debió
  renunciar por disposición del mismo Gómez) &
  \textbf{1928:} introduce el \textbf{Artículo 32, inciso G}: prohíbe las
  actividades con \textbf{propaganda comunista} y declara \textbf{traidor a la
  patria} a quien la haga. Vuelve a crear los cargos de \textbf{Presidente} y de
  \textbf{Comandante General del Ejército}. \\

  \textbf{1931--1935} & Gral. Gómez (4.ª vez; muere en la presidencia) &
  Instituye que el \textbf{presidente de la República sea a la vez Comandante
  General del Ejército}. \\
  \bottomrule
\end{longtable}
}

\begin{cuadroadvertencia}
  \small
  \textbf{Ojo con las fechas.} En la biografía, la fuente también habla de los
  períodos \textbf{1910--1915}, \textbf{1915--1922}, \textbf{1922--1929} y
  \textbf{1929--1936}. Es normal que las cronologías de un mismo período no
  coincidan exactamente: en un examen usa la secuencia que indique el profesor y
  \textbf{explica} el cambio de período con la reforma correspondiente.
\end{cuadroadvertencia}

\begin{cuadropasos}
  \small
  \textbf{Producto de la sesión.} Tabla de cuatro columnas en tu cuaderno:
  \textit{año -- norma -- qué cambió -- a quién benefició}.
\end{cuadropasos}

\subsection*{Preguntas de análisis}
\begin{preguntas}
  \item \nUno{} ¿En qué año el período presidencial pasó a \textbf{7 años} y se
        permitió la reelección?
  \item \nUno{} ¿Quién fue Juan Bautista Pérez y por qué debió renunciar?
  \item \nDos{} Explica la función del \textbf{presidente de la Alta Corte
        Federal y de Casación} en la reforma de 1914.
  \item \nTres{} Compara las reformas de \textbf{1909} y \textbf{1925}: ¿qué
        gana el Ejecutivo en cada una?
  \item \nTres{} ¿Por qué unas reformas \textbf{crean} el cargo de Comandante
        General del Ejército y otras lo \textbf{eliminan}? ¿Qué revela eso sobre
        dónde estaba el poder real?
  \item \nCuatro{} Detecta la discrepancia entre las dos cronologías que ofrece
        el texto. ¿Cómo la resolverías en un examen escrito sin contradecir al
        profesor?
\end{preguntas}

\pausa{Di en voz alta: \guillemotleft la reforma de \rule{1.6cm}{0.4pt} sirvió
para \rule{3cm}{0.4pt}\guillemotright.}

% ============================================================
\section{La fuerza: ejército, mano dura y nepotismo}
% ============================================================

\metasesion{25 + 5 min}{distinguir los tres pilares de coerción del régimen y
evaluar el costo social de cada uno.}

\begin{description}[leftmargin=1.2em, itemsep=5pt, style=nextline]
  \item[\textbf{Profesionalización militar.}]
  \textbf{Escuela Militar (1911)} y \textbf{Escuela de Aviación Militar (1920)}.
  \textbf{Ley del Servicio Militar Obligatorio (1926)} para contar con un
  ejército permanente.

  \item[\textbf{Control del territorio.}]
  Carreteras de \textbf{macadán} (algunas de asfalto) que unían ciudades y
  \textbf{plazas militares}.

  \item[\textbf{Desarme del pueblo.}]
  Desarmó a la población alegando que así \textbf{evitaba la delincuencia}.

  \item[\textbf{Represión.}]
  Exilio; cárceles: \textbf{La Rotunda} (Caracas), \textbf{Las Tres Torres}
  (Barquisimeto) y \textbf{Castillo de Puerto Cabello}; \textbf{trabajos
  forzados} en la construcción de carreteras y \textbf{torturas} frecuentes.

  \item[\textbf{Nepotismo.}]
  \textbf{Primera y segunda vicepresidencia} entregadas a su hermano
  \textbf{Juan Crisóstomo Gómez} (\guillemotleft Juancho\guillemotright) y a su
  hijo (\guillemotleft Vicentino\guillemotright).
\end{description}

\begin{cuadropasos}
  \small
  \textbf{Producto de la sesión.} Tres columnas ---\textit{coerción militar /
  coerción policial / coerción económica}--- con dos ejemplos en cada una.
\end{cuadropasos}

\subsection*{Preguntas de análisis}
\begin{preguntas}
  \item \nUno{} Menciona las dos escuelas militares creadas y el año de la Ley
        del Servicio Militar Obligatorio.
  \item \nUno{} ¿Qué tres cárceles se mencionan en la guía?
  \item \nDos{} ¿Qué relación existe entre \textbf{carreteras},
        \textbf{plazas militares} y control del territorio?
  \item \nTres{} ¿Qué función cumplían el \textbf{desarme del pueblo} y el
        argumento de \guillemotleft evitar la delincuencia\guillemotright?
        ¿A quién protegía realmente?
  \item \nCuatro{} ¿Por qué colocar a un hermano y a un hijo en las
        vicepresidencias es más que un privilegio económico? Relaciónalo con la
        \textbf{sucesión} del poder.
\end{preguntas}

\pausa{Nombra las tres cárceles sin mirar la guía.}

% ============================================================
\section{La economía I: la Venezuela agrícola (1908--1921)}
% ============================================================

\metasesion{25 + 5 min}{explicar cómo funcionaba la economía antes del petróleo
y quién se beneficiaba.}

\begin{itemize}[leftmargin=*, itemsep=3pt]
  \item El país seguía \textbf{dependiendo de los precios del café} en el mercado
        mundial: estos \textbf{cayeron}, mientras los del \textbf{cacao se
        mantuvieron}.
  \item Gómez \textbf{favoreció la agricultura} y buscó \textbf{acceso de los
        cultivadores al mercado internacional}.
  \item \textbf{Eliminó el monopolio de venta de tabaco}.
  \item Favoreció la \textbf{inversión extranjera} con facilidades para
        instalarse y \textbf{exención de impuestos} a los inversionistas.
  \item En \textbf{tenencia de tierras} fue el \textbf{mayor latifundista}:
        tenía tierras en \textbf{doce estados} del país y una parte estaban
        \textbf{ociosas}.
  \item La \textbf{industria ganadera} podía controlarse fácilmente: durante el
        período, la \textbf{venta de carne fue monopolizada} por el dictador y
        su camarilla. El General y sus amigos habían adquirido \textbf{tierras
        de cultivo y los mejores pastos}, \guillemotleft valiéndose de medios
        legales o de la fuerza\guillemotright.
\end{itemize}

\begin{cuadropasos}
  \small
  \textbf{Producto de la sesión.} Párrafo de cinco líneas:
  \guillemotleft La economía agrícola de Gómez benefició a
  \rule{1.8cm}{0.4pt} porque \rule{3.2cm}{0.4pt}\guillemotright.
\end{cuadropasos}

\subsection*{Preguntas de análisis}
\begin{preguntas}
  \item \nUno{} ¿De qué producto dependía la economía antes del petróleo y qué
        pasó con su precio?
  \item \nUno{} ¿Qué monopolio eliminó Gómez en la agricultura?
  \item \nDos{} ¿Qué tipo de facilidades ofrecía a los inversionistas
        extranjeros?
  \item \nTres{} ¿Cómo se relacionan \textbf{latifundio}, \textbf{monopolio de
        la carne} y \textbf{poder político}? Usa la frase
        \guillemotleft medios legales o de la fuerza\guillemotright.
  \item \nCuatro{} ¿La política agrícola de Gómez favoreció el
        \textbf{desarrollo nacional} o profundizó la \textbf{dependencia}? Toma
        posición con tres datos.
\end{preguntas}

\pausa{Repite la cadena: café $\rightarrow$ precios $\rightarrow$ exportaciones
$\rightarrow$ ingresos.}

% ============================================================
\section{La economía II: petróleo, crisis de 1929 y deuda}
% ============================================================

\metasesion{25 + 5 min}{explicar el cambio de modelo económico y evaluar la
paradoja de un país sin deuda pero dependiente.}

\subsection*{El modelo petrolero (desde 1921)}
\begin{itemize}[leftmargin=*, itemsep=3pt]
  \item Continúa la misma política de \textbf{ofertas y facilidades} al capital
        extranjero.
  \item \textbf{Derechos de exploración bajos}; \textbf{no} establece normas de
        \textbf{refinación ni perforación}.
  \item \textbf{Exenciones de impuestos de aduana} para maquinarias y equipos.
  \item \textbf{Plazos de concesión muy largos} y beneficiosos para las
        compañías.
  \item Concesiones a compañías \textbf{holandesas, norteamericanas e inglesas}.
  \item \textbf{1925:} el petróleo \textbf{desplaza al café} como primer producto
        de exportación.
  \item \textbf{1928:} Venezuela es \textbf{primer exportador mundial} de
        petróleo y \textbf{segundo productor}, después de Estados Unidos. Estados
        Unidos es el \textbf{primer inversionista} en el país.
\end{itemize}

\subsection*{La crisis económica de 1929}
\begin{center}
\small
\begin{tabularx}{0.85\linewidth}{@{}X l@{}}
  \toprule
  \textbf{Indicador} & \textbf{Efecto} \\
  \midrule
  Empleo petrolero             & Despidieron el \textbf{66\,\%} de los obreros \\
  Exportaciones petroleras     & Bajaron \\
  Ingresos fiscales            & Bajaron \\
  Exportaciones agrícolas      & Disminuyeron \\
  Importaciones                & Bajaron \\
  \bottomrule
\end{tabularx}
\end{center}

\subsection*{La deuda del país}
\begin{center}
\small
\begin{tabular}{@{}l r r@{}}
  \toprule
  \textbf{Año} & \textbf{Deuda externa} & \textbf{Deuda interna} \\
  \midrule
  1 de enero de 1909       & Bs.\ 161.138.109 & Bs.\ 71.958.383 \\
  1929                     & Bs.\ 29.249.416  & Bs.\ 28.735.884 \\
  1935 (muerte de Gómez)   & \multicolumn{2}{c}{\textbf{Saldada}} \\
  \bottomrule
\end{tabular}
\end{center}

\noindent Gómez se comportó como \guillemotleft el mejor pagador\guillemotright:
a su muerte \textbf{ya no había deuda}.

\begin{cuadropasos}
  \small
  \textbf{Producto de la sesión.} Gráfico simple de barras comparando la deuda
  total de 1909 (Bs.\ 233.096.492) con la de 1929 (Bs.\ 57.985.300).
\end{cuadropasos}

\subsection*{Preguntas de análisis}
\begin{preguntas}
  \item \nUno{} ¿En qué año el petróleo desplazó al café y en qué año Venezuela
        fue primer exportador mundial?
  \item \nUno{} Enumera \textbf{tres} efectos de la crisis de 1929 en Venezuela.
  \item \nDos{} Explica por qué los \textbf{plazos largos} y los
        \textbf{derechos de exploración bajos} beneficiaban a las compañías
        extranjeras.
  \item \nTres{} Analiza la aparente contradicción: un país que \textbf{paga
        toda su deuda} y, al mismo tiempo, \textbf{entrega concesiones} y no fija
        normas de refinación. ¿Cómo se sostiene ese equilibrio?
  \item \nTres{} ¿Qué muestra el dato del \textbf{66\,\% de obreros
        despedidos} sobre quién pagaba realmente la crisis?
  \item \nCuatro{} ¿Puede un país estar \textbf{sin deuda} y ser
        \textbf{dependiente} al mismo tiempo? Defiende tu respuesta con datos de
        esta sesión.
\end{preguntas}

\pausa{Escribe de memoria las cuatro cifras de la deuda.}

% ============================================================
\section{La sociedad: población, éxodo y estratos}
% ============================================================

\metasesion{25 + 5 min}{relacionar los cambios demográficos con la llegada del
petróleo.}

\subsection*{Población}
\begin{itemize}[leftmargin=*, itemsep=3pt]
  \item \textbf{Aumenta} de 1908 a 1911, \textbf{desciende} hasta 1920 y desde
        ahí \textbf{crece sin detenerse}.
  \item Causas del descenso (\textbf{1911--1920}): \textbf{mortalidad infantil},
        \textbf{desatención sanitaria}, la \textbf{epidemia de paludismo de 1916}
        y la \textbf{gripe española de 1918}.
  \item Causas del aumento (desde \textbf{1920}): inicio de la \textbf{época
        petrolera}, mejores condiciones sanitarias y \textbf{servicios de
        prevención y curación} instalados por las compañías petroleras para
        empleados y obreros.
\end{itemize}

\subsection*{Territorio}
\begin{itemize}[leftmargin=*, itemsep=3pt]
  \item Cambia la relación \textbf{población urbana / población rural}: se
        produce el \textbf{éxodo campesino}.
  \item El \textbf{área occidental} del país (donde empezó la exploración y
        explotación) y la ciudad de \textbf{Maracaibo} se hacen populosas.
\end{itemize}

\subsection*{Caudillismo}
\begin{itemize}[leftmargin=*, itemsep=3pt]
  \item \textbf{Desaparece el caudillismo regional}: el dictador es el gran
        caudillo. La propaganda lo presenta como \textbf{«Gómez Único»},
        \textbf{«el Rehabilitador»} y \textbf{«el Pacificador»}.
\end{itemize}

\subsection*{Estratos sociales}
\begin{enumerate}[leftmargin=*, itemsep=2pt]
  \item \textbf{Terratenientes} (con Gómez, sus familiares y amigos).
  \item \textbf{Comerciantes, banqueros e industriales}, nacionales y
        extranjeros.
  \item \textbf{Burocracia} de altos empleados gubernamentales.
  \item \textbf{Empleados y profesionales} en general.
  \item \textbf{Obreros}.
  \item \textbf{Campesinos y peones} de haciendas y hatos.
\end{enumerate}

\begin{cuadropasos}
  \small
  \textbf{Producto de la sesión.} Pirámide social dibujada con los seis estratos
  y una palabra que describa la situación de cada uno.
\end{cuadropasos}

\subsection*{Preguntas de análisis}
\begin{preguntas}
  \item \nUno{} ¿Qué epidemias se mencionan y en qué años?
  \item \nUno{} ¿Qué dos zonas del país se hicieron populosas con el petróleo?
  \item \nDos{} Explica con tus palabras qué es el \textbf{éxodo campesino} y por
        qué ocurre.
  \item \nTres{} ¿Qué relación hay entre \textbf{crecimiento demográfico},
        \textbf{petróleo} y \textbf{servicios de salud}?
  \item \nTres{} ¿Por qué la desaparición del caudillismo regional puede
        describirse como \textbf{concentración} y no como \textbf{fin} del
        caudillismo?
  \item \nCuatro{} ¿Quiénes \textbf{no} aparecen en la lista de estratos con
        poder? ¿Qué te dice esa ausencia sobre la sociedad gomecista?
\end{preguntas}

\pausa{Ordena los seis estratos de arriba a abajo sin mirar.}

% ============================================================
\section{La cultura: educación, positivismo y prensa}
% ============================================================

\metasesion{25 + 5 min}{explicar la contradicción entre modernización material y
atraso educativo.}

\subsection*{Educación}
\begin{itemize}[leftmargin=*, itemsep=3pt]
  \item Estuvo \textbf{desasistida} durante el gomecismo: la \textbf{Universidad
        Central permaneció cerrada desde 1912 hasta 1925}.
  \item El \textbf{analfabetismo adulto} se situaba entre el \textbf{70\,\%} y
        el \textbf{90\,\%}: la fuente indica que hasta el 90\,\% de la población
        no sabía leer ni escribir.
  \item Los \textbf{presupuestos anuales} del Ministerio de Educación eran
        \textbf{ínfimos}.
  \item La \textbf{inscripción escolar en primaria} fue considerablemente
        \textbf{inferior} a la de la época guzmancista: en \textbf{1883} había
        \textbf{91.462 alumnos} inscritos, un \textbf{4,45\,\%} de la población
        total.
\end{itemize}

\subsection*{Pensamiento}
\begin{itemize}[leftmargin=*, itemsep=3pt]
  \item Continúa el movimiento histórico-sociológico del \textbf{positivismo},
        presente desde Castro, con \textbf{José Gil Fortoul}, \textbf{Laureano
        Vallenilla Lanz}, \textbf{César Funesta}, \textbf{Pedro Manuel Arcaya} y
        \textbf{Lisandro Alvarado}.
\end{itemize}

\subsection*{Literatura y prensa}
\begin{itemize}[leftmargin=*, itemsep=3pt]
  \item La \textbf{novela} tiene en \textbf{Rómulo Gallegos} a su máximo
        representante.
  \item Revistas importantes: \textbf{Válvula} y \textbf{El Cojo Ilustrado}
        (esta última dura hasta \textbf{1915}).
  \item \textbf{Periodismo humorístico} como vía de crítica indirecta:
        \textbf{Fantoches}, de \textbf{Leoncio Martínez} («Leo»), y
        \textbf{Pitorreos}, de \textbf{Francisco Pimentel} («Job Pim»).
  \item Prensa oficial: \textbf{El Nuevo Diario}. En el interior,
        \textbf{El Impulso} de Barquisimeto era uno de los diarios más
        importantes.
\end{itemize}

\begin{cuadropasos}
  \small
  \textbf{Producto de la sesión.} Cuadro de dos columnas:
  \textit{lo que modernizó Gómez / lo que dejó atrasado}.
\end{cuadropasos}

\subsection*{Preguntas de análisis}
\begin{preguntas}
  \item \nUno{} ¿Cuántos años estuvo cerrada la Universidad Central?
  \item \nUno{} Menciona \textbf{tres} representantes del positivismo
        venezolano.
  \item \nDos{} ¿Qué papel cumplía el \textbf{periodismo humorístico} en un
        régimen sin partidos ni libertad de prensa?
  \item \nTres{} Relaciona el dato de \textbf{1883} (91.462 alumnos =
        4,45\,\%) con la idea de una educación \guillemotleft
        desasistida\guillemotright. ¿Qué te dice la comparación con la época
        anterior?
  \item \nTres{} ¿Por qué el cierre de la Universidad Central es un dato
        \textbf{político} y no solo educativo?
  \item \nCuatro{} Gómez modernizó el ejército y las carreteras, pero no la
        escuela. Formula una \textbf{hipótesis} que explique esa elección y
        respáldala con dos datos de la guía.
\end{preguntas}

\pausa{Escribe el nombre de las dos revistas y de los dos periódicos
humorísticos.}

% ============================================================
\section{La oposición y la Generación del 28}
% ============================================================

\metasesion{25 + 5 min}{comparar las tres corrientes opositoras y explicar por
qué solo una tuvo proyección histórica.}

La oposición tuvo \textbf{tres matices} según sus protagonistas.

\subsection*{a) Los antiguos caudillos}
\begin{itemize}[leftmargin=*, itemsep=3pt]
  \item Posición \textbf{aislada, sin fuerza y sin consenso}; varios líderes
        insistentes en invasiones y alzamientos.
  \item Fracasan por la \textbf{cohesión, preparación y equipamiento} del
        ejército del dictador.
  \item Terminan en el \textbf{exilio}, fusilados, en la cárcel o torturados.
\end{itemize}

{\small
\begin{tabular}{@{}c >{\raggedright\arraybackslash}p{2.7cm}
                    >{\raggedright\arraybackslash}p{5.5cm}
                    >{\raggedright\arraybackslash}p{6.3cm}@{}}
  \toprule
  \textbf{Año} & \textbf{Caudillo} & \textbf{Acción} & \textbf{Resultado} \\
  \midrule
  1913 & Gral. Román Chalbaud &
  Gran conspiración en Caracas &
  Fracasa; encerrado en La Rotunda hasta 1927, luego exiliado \\
  1914 & Gral. Emilio Arévalo Cedeño &
  Invasión por la frontera de Colombia (Táchira) &
  Fracasa; repite en 1915 y 1920 con el mismo resultado \\
  1918 & Gral. Juan Pablo Peñaloza &
  Invasión por el Táchira &
  Fracasa \\
  1927 & Gral. Rafael Simón Urbina &
  Invasión por Coro &
  Fracasa \\
  \bottomrule
\end{tabular}
}

\subsection*{b) Los oficiales jóvenes}
Conspiraciones dentro de los cuarteles de Caracas:
\begin{itemize}[leftmargin=*, itemsep=3pt]
  \item \textbf{1919:} figura principal \textbf{Luis Rafael Pimentel}.
  \item \textbf{1922:} participan, entre otros, \textbf{Alejandro Rascaniere} y
        \textbf{Francisco Angarita Arvelo}.
  \item En ambos casos los conspiradores terminaron en la \textbf{cárcel}.
\end{itemize}

\subsection*{c) Los estudiantes: la Generación del 28}
\begin{itemize}[leftmargin=*, itemsep=3pt]
  \item Grupo de universitarios que protagonizaron en el \textbf{carnaval
        caraqueño de 1928} un movimiento \textbf{académico y estudiantil},
        iniciado con la \textbf{semana del estudiante} y que culminó en un
        \textbf{enfrentamiento con el régimen}.
  \item Lo que empezó como un proyecto \textbf{restringido a la Universidad} se
        transformó en una propuesta de \textbf{cambio del régimen político} y de
        los \textbf{fundamentos de la sociedad y la cultura} venezolanas.
  \item Participaron, entre otros: \textbf{Rómulo Betancourt}, \textbf{Jóvito
        Villalba}, \textbf{Andrés Eloy Blanco} y \textbf{Raúl Leoni}.
  \item De ella salieron los fundadores de \textbf{Acción Democrática (AD)},
        \textbf{Unión Republicana Democrática (URD)} y el \textbf{Partido
        Comunista}.
\end{itemize}

\begin{cuadropasos}
  \small
  \textbf{Producto de la sesión.} Cuadro comparativo de las tres corrientes:
  \textit{¿quiénes? / ¿cómo actuaron? / ¿por qué fracasaron o triunfaron?}
\end{cuadropasos}

\subsection*{Preguntas de análisis}
\begin{preguntas}
  \item \nUno{} ¿Cuáles son las tres corrientes de oposición al gomecismo?
  \item \nUno{} Menciona \textbf{dos} invasiones fallidas con año y caudillo.
  \item \nDos{} ¿Por qué fracasaban los alzamientos caudillistas?
  \item \nTres{} ¿Qué tuvo de \textbf{distinto} la Generación del 28 frente a los
        caudillos y frente a los oficiales jóvenes?
  \item \nTres{} ¿Por qué un movimiento que comenzó como \guillemotleft semana
        del estudiante\guillemotright{} pudo convertirse en un proyecto político
        nacional?
  \item \nCuatro{} Explica cómo la Generación del 28 conecta con el nacimiento de
        \textbf{AD}, \textbf{URD} y el \textbf{PCV}. ¿Qué continuidad hay entre
        la protesta estudiantil y los partidos del siglo XX?
\end{preguntas}

\pausa{Nombra los cuatro estudiantes citados.}

% ============================================================
\section{Cierre: síntesis, cronología y simulacro}
% ============================================================

\metasesion{25 + 5 min}{integrar todo el objetivo 2.1 en una sola respuesta
escrita.}

\subsection*{1. Cronología integradora (complétala sin mirar)}

{\small
\begin{tabular}{@{}c p{14.2cm}@{}}
  \toprule
  \textbf{Año} & \textbf{Hecho} \\
  \midrule
  1857 & \dotfill \\
  1883 & \dotfill \\
  1899 & \dotfill \\
  1902 & \dotfill \\
  1908 & \dotfill \\
  1911 & \dotfill \\
  1912--1925 & \dotfill \\
  1916 y 1918 & \dotfill \\
  1920 & \dotfill \\
  1921 & \dotfill \\
  1925 & \dotfill \\
  1926 & \dotfill \\
  1928 & \dotfill \\
  1929 & \dotfill \\
  1930 & \dotfill \\
  1931 & \dotfill \\
  1935 & \dotfill \\
  \bottomrule
\end{tabular}
}

\subsection*{2. Glosario mínimo}
\noindent\small
\textbf{Gomecismo} · \textbf{caudillismo} · \textbf{fuerzas vivas} ·
\textbf{nepotismo} · \textbf{latifundio} · \textbf{concesión} ·
\textbf{macadán} · \textbf{La Rotunda} · \textbf{positivismo} ·
\textbf{Generación del 28} · \textbf{co-presidencia} ·
\textbf{propaganda comunista (Art.\ 32, inciso G)}.

\subsection*{3. Preguntas integradoras}
\begin{preguntas}
  \item \nCuatro{} ¿Cómo se relacionan \textbf{reforma constitucional},
        \textbf{ejército} y \textbf{petróleo} para sostener 27 años de
        dictadura? Responde en un párrafo con las tres ideas conectadas.
  \item \nCuatro{} La guía muestra un país que \textbf{moderniza} carreteras y
        ejército, \textbf{paga} su deuda y, a la vez, \textbf{cierra} la
        universidad y \textbf{monopoliza} la carne. ¿Qué tipo de modernización
        es esa? Defínela con tus palabras.
  \item \nCuatro{} ¿Por qué la \textbf{oposición armada} fracasó y la
        \textbf{oposición estudiantil} tuvo proyección? Compara recursos,
        organización y discurso.
  \item \nCuatro{} ¿Qué cambió y qué permaneció entre \textbf{1908} y
        \textbf{1935} en la economía venezolana? Usa la idea de
        \textit{transición} (agrícola $\rightarrow$ petrolera) sin simplificar.
  \item \nCuatro{} Evalúa el legado: ¿Gómez dejó un \textbf{Estado} más fuerte o
        un \textbf{país} más débil? Argumenta con dos datos económicos y dos
        sociales.
\end{preguntas}

\subsection*{4. Autoevaluación}
\noindent\small
Marca \textbf{\checkmark} solo lo que puedas hacer \textbf{sin mirar la guía}:

\begin{itemize}[leftmargin=*, itemsep=2pt, label=$\square$]
  \item Explico por qué Gómez llega al poder en 1908.
  \item Ordeno las reformas constitucionales y digo qué cambió en cada una.
  \item Diferencio la Venezuela agrícola de la petrolera con tres datos cada una.
  \item Explico la crisis de 1929 con cuatro efectos.
  \item Describo los seis estratos sociales.
  \item Explico la contradicción entre modernización material y atraso educativo.
  \item Comparo las tres corrientes de oposición.
  \item Escribo en 10 líneas una respuesta completa al objetivo 2.1.
\end{itemize}

\begin{cuadroextra}
  \small
  \textbf{Cierre del objetivo.} Si marcaste menos de 6 casillas, repite las
  sesiones de los temas sin marcar. Si marcaste 6 o más, estás listo para el
  examen.
\end{cuadroextra}

% ============================================================
\subsection*{Fuentes citadas por la guía original}
\noindent\small
Rena (Red Educativa Nacional) · Venezuela Tuya · Monografías.com · Ensayos en
línea (Espacio IsaurajRiver).

\vspace{6pt}
\begin{cuadroinfo}
  \footnotesize
  \textbf{Nota final.} Esta versión reordena, titula y añade preguntas;
  \textbf{no} sustituye la explicación del Prof.\ José Rambert ni sus fuentes.
  Reedición elaborada con su autorización.
\end{cuadroinfo}

\end{document}
